\section{Análisis de paquetes capturados: comprensión a nivel de la API}

El expositor realizó un análisis de paquetes capturados de la API de Gemini mediante un hook de solicitudes/respuestas HTTP:

\begin{itemize}
    \item \textbf{Request Body}: contiene la definición de las Tools (Code Execution) y la entrada del usuario (imagen + texto)
    \item \textbf{Response Body}: contiene varias Parts, cada una de las cuales puede ser:
    \begin{itemize}
        \item \texttt{executable\_code}: código Python generado por el modelo
        \item \texttt{code\_execution\_result}: resultado de la ejecución del código (incluidas las imágenes generadas)
        \item \texttt{text}: salida de razonamiento textual del modelo
    \end{itemize}
    \item cada paso del proceso que muestra la interfaz es, en esencia, la visualización de estas Parts
\end{itemize}

\begin{knowledgebox}{Gemini App vs AI Studio}
Al momento de grabarse el video, la Gemini App (la versión para consumidores) aún no ofrecía la Code Execution Tool. Para experimentar Agentic Vision es necesario usar AI Studio (la herramienta para desarrolladores).
\end{knowledgebox}

